\chapter{If it happens}

Two clocks run at once. The human clock is first. The
chain clock is second, and only if speeding it up does
not make the first one worse.

\section{Limits that buy time}

\textbf{Hard caps} on hot and warm wallets so a coerced
session cannot drain the reserve.

\textbf{Extra approval} once a transfer crosses a line,
even if the request looks like payroll. Send that
approval on a different channel, to a different person,
on a different device. One compromised chat should not
complete the payment.

\textbf{Delay} on large moves. The attacker has to keep
control. That is when someone else can notice.

Request and approval should not live on the same phone.
High-privilege signers should not be in the backpack you
take to the airport. Geographic spread means a local
incident cannot collect every signature that night.

\section{A path that can be complied with}

Keep a low-value path that can satisfy an attacker
quickly and does not upgrade into cold. It has to match
the target's public life well enough to survive a glance
at a block explorer. It must not be a stepping stone.

During travel, lower the limits and raise the approvals
on purpose. Credible threats, doxxing, or repeated
unfamiliar surveillance are the same trigger: shrink what
the org can do until you understand the situation.

\section{Life first}

Treat it as a life-safety emergency. Assume the victim
is still being watched or still under threat.

Activate one crisis channel and one incident commander.
Split brains make this worse. If kidnap-and-ransom
insurance exists, call that team. Reputable policies
come with negotiators. Do not crowdsource this in a
group chat.

Try to learn what the attacker likely already has:
phones, seeds, which signer, which room.

\textbf{Do not instruct resistance} while someone is
being hurt. Move people to safety, get medical care, then
debrief. If this started as travel, assume device seizure
or detention and get local legal help before you play
incident-responder on Telegram.

\section{Containment when it is safe}

When the violence is no longer the current problem:

Rotate or disable signers. Move what is still exposed
into a contingency path you already built.

If contracts can pause, timelock, or restrict admin
functions, consider it. Coercion is a way to pass an
upgrade, not only a way to steal a hot wallet.

Assume every secret on an unlocked device is gone.
Rebuild on clean hardware. Do not ``just change the PIN.''

\section{Other people}

Custodians, exchanges, vendors: use the escalation path
you wrote down beforehand. Block movements. Write down
what you asked them to do.

Law enforcement through counsel and a crisis team, in a
way that does not publish the victim's location or the
remaining signers.

Insurance notice per the policy, without dumping signer
maps into an email that will be forwarded.

\section{Write it down, show it to few}

Keep hashes, times, messages, device state. Keep a log of
who decided what. Both are for later. Neither is for the
company all-hands.

Need-to-know is a safety control here. Extra readers
become extra targets.

When the immediate incident is over, the next chapter
starts. Do not skip the people and jump to a new seed
ceremony because it feels like action.
